\section{Quadratic feedback: finite coefficients and boundary inversion}
\label{qfb:section}

Order jets describe how a polylogarithm changes with its order. Nonlinear
feedback asks a complementary question: what remains explicit when the
unknown also occurs in the argument? Consider the germ
\begin{equation}\label{qfb:model}
 uv(r)-\frac{v(r)^2}{2}=\Li_s(re^{v(r)}),\qquad v(0)=0,\qquad u>0,
 \qquad a_n(s,u)=[r^n]v(r).
\end{equation}
The feedback reports \cite{qfb:reports} establish its interior inverse
geometry and complete the boundary regimes. We first develop an exact
finite coefficient formula, exposing a classical Bessel-polynomial
structure. We then reconcile the global radius theorem. This order keeps
finite identities, local expansions and global continuation distinct.

\subsection{Two Lagrange inversions and a Bessel-polynomial formula}
Let
\[
 C_{n,k}(s)=
 \sum_{\substack{m_1+\cdots+m_k=n\\m_i\ge1}}(m_1\cdots m_k)^{-s}
 =[t^n]\Li_s(t)^k .
\]
The sum is finite, so it is entire in $s$. Define the classical Bessel
polynomial \cite{qfb:bessel}
\[
 y_d(z)=\sum_{j=0}^d\frac{(d+j)!}{(d-j)!\,j!}\left(\frac z2\right)^j,
 \qquad y_0=1,\quad y_1=1+z.
\]
The positivity of its coefficients will be useful, but its role here is
an exact change of coordinates, not an asymptotic approximation.

\begin{theorem}[Finite Bessel coefficient identity]\label{qfb:bessel-theorem}
For $n\ge1$, $s\in\mathbb C$ and $u>0$,
\begin{align}
 a_n(s,u)&=\sum_{k=1}^n Q_{n,k}(u)C_{n,k}(s),\label{qfb:finite}\\
 Q_{n,k}(u)&=\frac{n^{k-1}}{k!\,u^k}
                 y_{k-1}\left(\frac1{nu}\right)\notag\\
 &=\sum_{j=0}^{k-1}
 \frac{(k-1+j)!\,n^{k-1-j}}
 {k!\,j!\,(k-1-j)!\,2^j u^{k+j}}.\label{qfb:weight}
\end{align}
More generally, replace $\Li_s(t)$ by any holomorphic germ
$A(t)$ with $A(0)=0$ and replace $C_{n,k}$ by $[t^n]A(t)^k$.
The same formula holds for every complex $u\ne0$.
\end{theorem}
\begin{proof}
The implicit derivative of the defining equation in $v$ at the origin
is $u\ne0$, so the germ exists and is unique. Set
$f(t)=u-\sqrt{u^2-2A(t)}$, choosing the square root with value $u$.
Then $v=f(t)$ and $t=re^{f(t)}$. Lagrange inversion gives
\[
 [r^n]v=\frac1n[t^{n-1}]f'(t)e^{nf(t)}
       =\frac1n[t^n]e^{nf(t)}.
\]
For the second inversion write $g(z)=u-\sqrt{u^2-2z}$.
Its inverse is $z=g(u-g/2)$, whence
\[
 [z^k]e^{ng(z)}
 =\frac nk[w^{k-1}]e^{nw}(u-w/2)^{-k}.
\]
Expand both factors on the right and divide by the first factor $n$.
The result is \eqref{qfb:weight}. Composing $z=A(t)$ requires only
$1\le k\le n$, because $A(0)=0$. This proves both assertions.
\end{proof}

For fixed $n$, the weights have a short exact recurrence:
\begin{equation}\label{qfb:weight-recurrence}
 k(k+1)u^2Q_{n,k+1}
 =k(2k-1)Q_{n,k}+n^2Q_{n,k-1}\qquad(k\ge2),
\end{equation}
with $Q_{n,1}=u^{-1}$ and
$Q_{n,2}=n/(2u^2)+1/(2u^3)$. It follows either by comparing
\eqref{qfb:weight}, or from
$y_k=(2k-1)zy_{k-1}+y_{k-2}$.
For example,
\begin{align*}
 a_1&=\frac1u,\\
 a_2&=\frac{2^{-s}}u+\frac1{u^2}+\frac1{2u^3},\\
 a_3&=\frac{3^{-s}}u+
       2^{-s}\left(\frac3{u^2}+\frac1{u^3}\right)
       +\frac3{2u^3}+\frac3{2u^4}+\frac1{2u^5}.
\end{align*}
For real $s$, $a_n$ is a polynomial in $u^{-1}$ with strictly positive
coefficients at every degree from $1$ through $2n-1$. The highest
coefficient is $\operatorname{Cat}_{n-1}/2^{n-1}$, where
$\operatorname{Cat}_d=(2d)!/(d!(d+1)!)$; for $n\ge2$ the next is
$n\operatorname{Cat}_{n-1}/2^{n-1}$. Only the composition into $n$
ones contributes to these last two powers, so they are independent of $s$.

\subsection{Finite order spectra and exact recurrences}
Let $\mathcal M_n$ be the set of products of all ordered compositions
of $n$, with arbitrary positive length. If $N_{n,k}(M)$ counts the
compositions of length $k$ and product $M$, put
\[
 W_{n,M}(u)=\sum_{k=1}^n Q_{n,k}(u)N_{n,k}(M).
\]
\begin{corollary}[A finite Dirichlet polynomial in the order]
\label{qfb:order-spectrum}
For $u>0$,
\begin{equation}\label{qfb:dirichlet}
 a_n(s,u)=\sum_{M\in\mathcal M_n}W_{n,M}(u)M^{-s},
 \qquad W_{n,M}(u)>0.
\end{equation}
Consequently
\begin{equation}\label{qfb:annihilator}
 \prod_{M\in\mathcal M_n}(\partial_s+\log M)\,a_n(s,u)=0.
\end{equation}
This is the monic constant-coefficient annihilator of least degree.
For a real $h\ne0$, define
$P_{n,h}(X)=\prod_{M\in\mathcal M_n}(X-M^{-h})
 =\sum_{j=0}^d p_jX^j$. Then, identically in $s$,
\begin{equation}\label{qfb:order-shift}
 \sum_{j=0}^d p_j a_n(s+jh,u)=0.
\end{equation}
In particular $W_{n,1}=Q_{n,n}$ and
$\lim_{s\to+\infty}a_n(s,u)=Q_{n,n}(u)$.
\end{corollary}
\begin{proof}
Group the finite sum \eqref{qfb:finite} by the composition product.
Every product in the support occurs with a positive weight. Applying a
constant-coefficient differential operator with polynomial $P$ gives
$\sum_M W_{n,M}P(-\log M)M^{-s}$. Distinct real exponentials are
linearly independent: differentiating at one point gives a Vandermonde
matrix in their distinct exponents. Hence any annihilator must vanish
at every $-\log M$, proving minimality. The shift formula follows by
the same calculation with the distinct eigenvalues $M^{-h}$.
The product is one only for the composition entirely of ones, proving
the final statements.
\end{proof}

This finite spectrum concerns a fixed Taylor coefficient, not the
complete function as $n$ varies. All order derivatives are explicit:
$\partial_s^j a_n=\sum_M W_{n,M}(-\log M)^jM^{-s}$.
For instance $\mathcal M_3=\{1,2,3\}$, so the unit-order samples satisfy
\[
 a_3(s+3,u)-\frac{11}{6}a_3(s+2,u)
          +a_3(s+1,u)-\frac16a_3(s,u)=0.
\]
This exact relation survives every choice of positive feedback $u$.

\subsection{The positive branch and its exact radius}
The finite formula proves $a_n>0$ for real $s$. Now suppose $s>0$
and write, for $s>2$,
\[
 A=\zeta(s-1),\qquad U_c(s)=\sqrt{2\zeta(s)+A^2}.
\]
\begin{theorem}[Complete quadratic-feedback classification]
\label{qfb:phase}
The radius $R$ of \eqref{qfb:model} is the unique singular point on its
circle of convergence, and the Taylor series converges absolutely there.
Its value and coefficient equivalents are as follows.

If $0<s\le2$, or if $s>2$ and $u<U_c(s)$, let $t_*\in(0,1)$ be
the unique solution of
\[
 u^2=2\Li_s(t_*)+\Li_{s-1}(t_*)^2,\qquad w_*=\Li_{s-1}(t_*).
\]
Then
\begin{equation}\label{qfb:interior}
 R=t_*e^{w_*-u},\qquad
 a_n\sim\sqrt{\frac{w_*}{2\pi(1+\Li_{s-2}(t_*))}}\,
                  R^{-n}n^{-3/2}.
\end{equation}
For $s>2$ and $u\ge U_c(s)$ set
$b=\sqrt{u^2-2\zeta(s)}$, $v_c=u-b$, $R_b=e^{-v_c}$.
Here $R=R_b$. If $b>A$, then
\begin{equation}\label{qfb:supercritical}
 a_n\sim\frac{b^{s-1}}{(b-A)^s}R_b^{-n}n^{-s}.
\end{equation}
If $b=A$, the three critical laws are
\begin{equation}\label{qfb:critical}
 a_n\sim R_b^{-n}
 \begin{cases}
 \displaystyle
 \frac{(A/\Gamma(1-s))^{1/(s-1)}}{-\Gamma(-1/(s-1))}
                 n^{-1-1/(s-1)},&2<s<3,\\[6pt]
 \displaystyle \sqrt{A/\pi}\,n^{-3/2}(\log n)^{-1/2},&s=3,\\[4pt]
 \displaystyle \sqrt{\frac{A}{2\pi(1+\zeta(s-2))}}\,n^{-3/2},&s>3.
 \end{cases}
\end{equation}
\end{theorem}
\begin{proof}
The parametrization from the first Lagrange inversion has
\[
 \frac{d}{dt}\log r(t)=
 \frac1t\left(1-\frac{\Li_{s-1}(t)}
                    {\sqrt{u^2-2\Li_s(t)}}\right).
\]
The characteristic right side $2\Li_s(t)+\Li_{s-1}(t)^2$ increases
strictly from zero, with endpoint infinity for $s\le2$ and $U_c(s)^2$
otherwise. Any characteristic point precedes a zero of the radicand.
Thus the positive branch increases to the stated fold or endpoint.

For a global convergence argument, iterate
$\Phi(r,v)=v^2/(2u)+u^{-1}\sum_{m\ge1}r^me^{mv}/m^s$ from zero.
The coefficients are nonnegative and stabilize degree by degree; for
each positive $r$ below the proposed radius, the real iterates are
bounded by the parametrized solution. The Taylor series therefore
converges there and agrees with that solution by local uniqueness.
Monotone convergence at the endpoint shows $\sum_n a_nR^n=v_c$,
or $u-w_*$ in the interior case.

At $|r|=R$, $r\ne R$, strict positivity in every degree implies
$|v(r)|<v(R)$. Hence $|re^{v(r)}|<1$, or $<t_*$ in the interior
case. Moreover
\[
 |\Phi_v(r,v(r))|
 \le\frac{|v(r)|+\Li_{s-1}(|re^{v(r)}|)}u<1.
\]
At the positive endpoint the upper comparison is
$(v_c+A)/u\le1$, or $(u-w_*+w_*)/u=1$.
The implicit function theorem extends across every other point of the
circle. It remains to prove that the positive endpoint is singular.

For the boundary cases put $x=-\operatorname{Log}t$,
$d=-\operatorname{Log}(r/R_b)$ and
\[
 D_s(x)=\zeta(s)-\Li_s(e^{-x}),\qquad
 K_s(x)=Ax+\frac{x^2}{2}-D_s(x).
\]
The selected square root gives the exact identities
\begin{align}
 (b+x)d-\frac{d^2}{2}&=(b-A)x+K_s(x),&
 v&=v_c+d-x,\label{qfb:normal}\\
 K_s(x)&=\int_0^x(x-t)\{1+\Li_{s-2}(e^{-t})\}\,dt.\notag
\end{align}
The integral converges at zero for $s>2$. The paired expansion
\eqref{rj:lerch} at $a=1$ gives, for noninteger $s$,
\[
 K_s(x)=\Gamma(1-s)x^{s-1}
       +\frac{1+\zeta(s-2)}2x^2+O(x^3).
\]
At $s=3$ it gives instead
\begin{equation}\label{qfb:K3}
 K_3(x)=\frac{x^2}{2}\left(\log\frac1x+\frac52\right)
           +\frac{x^3}{12}-\frac{x^4}{288}+O(x^6).
\end{equation}
The $5/2$ includes the quadratic feedback contribution.

If $b>A$, put $\kappa=1-A/b>0$. The first nonanalytic term
of $v$ after inversion is $\Gamma(1-s)d^{s-1}/(b\kappa^s)$.
At integer $s=m\ge3$, pairing first gives the term
$(-1)^m d^{m-1}\log d/(b(m-1)!\kappa^m)$.
Power or integer-logarithm transfer cancels the Gamma factor or
the signed factorial and proves \eqref{qfb:supercritical}.

At $b=A$, the leading terms of $d$ are
$\Gamma(1-s)x^{s-1}/A$, $x^2\log(1/x)/(2A)$, or
$(1+\zeta(s-2))x^2/(2A)$ in the three regimes.
Inverting and using $v=v_c+d-x$ gives \eqref{qfb:critical} by
power--logarithm transfer \cite{qfb:transfer}. In particular the middle
inverse has leading term $2\sqrt A\,d^{1/2}(\log(1/d))^{-1/2}$.
At an interior fold, expansion of the parametrization gives the
square-root amplitude $\sqrt{2w_* /(1+\Li_{s-2}(t_*))}$,
which proves \eqref{qfb:interior}.

For completeness these transfers have the required complex domains.
For $d\sim cx^p$ on a closed slit sector, use the model inverse
$(d/c)^{1/p}$; for $d\sim cx^2\log(1/x)$ use
$(2d/(c\log(1/d)))^{1/2}$. On a fixed small relative circle around
that model, the normalized convergent expansion tends uniformly to
the corresponding power equation. Rouch\'e's theorem gives one root.
Differentiated convergent expansions make its derivative nonzero,
so the local inverses glue on smaller slit sectors.
The analytic coordinate $d=1-r/R_b+O((1-r/R_b)^2)$ supplies the
local dent. A finite cover of the remaining circle, using the already
proved implicit continuations, completes the transfer domain.
The nonanalytic leading terms prove genuine singularity and the
asserted radius.
\end{proof}

\begin{corollary}[Physical evaluation threshold]\label{qfb:physical}
For each $s>0$, let $t_s\in(0,1)$ be the unique solution of
$\log t_s+\Li_{s-1}(t_s)=0$, and put
$u_s^2=2\Li_s(t_s)+\Li_{s-1}(t_s)^2$.
The Taylor series at $r=-e^{-u}$ converges exactly when $u\ge u_s$,
and converges absolutely at equality.
\end{corollary}
\begin{proof}
The left side increases strictly from negative infinity to a positive
limit. In the interior regime $e^{-u}/R=e^{-w_*}/t_*$, giving exactly
the stated threshold. For $s>2$ it occurs below $U_c(s)$; at and above
that transition $e^{-u}/R_b=e^{-b}<1$. Outside the radius the series
diverges, and the theorem gives absolute convergence on its circle.
\end{proof}

\subsection{Evidence and the next inverse questions}
The interior result and boundary completion retain the attribution in
\cite{qfb:reports}. The latter resolves an explicitly posed project
question; it is not a correction to a theorem that claimed those cases.
The finite Bessel formula and order-spectrum identities above develop
the exact Lagrange representation using classical polynomial machinery.
They make no publication-priority claim.

Independent rational computations compare \eqref{qfb:finite} with
coefficient iteration of the original equation, including signed input
kernels; separate checks test the weight and order-shift recurrences.
The incoming full suites also reproduce the marginal coefficient ratio
$0.739119713$ to its leading estimate at $n=2048$, versus $1.010075657$
after the second logarithmic refinement. The finite discrepancy motivates
the refined inverse charts but is not a bound or a proof of an equivalent.

The report's convergent critical charts, complete logarithmic correction
calculus and joint crossover refinements retain separate integration
scopes. Its local complex folds do not settle global access from the
Taylor germ. Also, when $b<A$ is real, $R_b=e^{b-u}$ is an endpoint
reference smaller than the actual interior radius. The compound-Poisson
coefficient identity uses that reference on both sides of criticality;
it must not silently replace the radius established above.
